\documentclass[12pt]{article}

% ---------- Page Layout (Screenplay Standard) ----------
\usepackage[letterpaper,left=1.5in,right=1in,top=1in,bottom=0.9in]{geometry}

% ---------- Font & Encoding ----------
\usepackage[T1]{fontenc}
\usepackage[utf8]{inputenc}
\usepackage{courier} % Screenplay font
\renewcommand{\familydefault}{\ttdefault}

% ---------- Spacing ----------
\usepackage{setspace}
\setstretch{1.0} % Screenplay single spacing
\setlength{\parindent}{0pt}
\setlength{\parskip}{0pt}

% ---------- Microtypography ----------
\usepackage{microtype}

% ---------- Section Formatting (Scene-Oriented) ----------
\usepackage{titlesec}

% Scene headings (used for SCENE I, SCENE II, etc.)
\titleformat{\section}
  {\normalfont\bfseries}
  {}
  {0pt}
  {}
\titlespacing*{\section}{0pt}{2.5ex}{1.5ex}

% Subsections (optional: beats, inserts, interludes)
\titleformat{\subsection}
  {\normalfont\bfseries}
  {}
  {0pt}
  {}
\titlespacing*{\subsection}{0pt}{2.0ex}{1.0ex}

% ---------- Screenplay Line Breaking Fix ----------
\raggedright
\sloppy
\setlength{\emergencystretch}{3em}

\begin{document}

% ---------- Title Page ----------
\thispagestyle{empty}

\begin{center}
\vspace*{2.5cm}

{\Large\bfseries FLOWER WARS}\\[1.2cm]

{\normalsize A Screenplay}\\[3.5cm]

{\normalsize
Flyxion\\
\today\\
Draft 02
}

\vfill
\end{center}

\newpage
\section*{Log line}

In a rising Mexica empire strained by endless conquest, Tlacaelel—statesman, strategist, and reluctant reformer—reimagines war itself, transforming annihilation into ritualized conflict in an attempt to preserve cosmic order, political stability, and human life, even as the limits of restraint are tested by enemies who refuse to recognize the field.

\section*{Scene I-A}

\textbf{INT. TLACAELEL’S FAMILY HOME — NIGHT}

Firelight moves softly across plastered walls painted with red and black geometric patterns. The house is modest by noble standards but carefully kept. A single room opens onto a small inner courtyard where night air carries distant drums and the faint smell of lake water.

A bark-paper \textbf{CODEX} lies open on the floor. Its pages show tribute tallies, campaign routes, and genealogical glyphs. Small \textbf{OBSIDIAN COUNTING STONES} are arranged in deliberate rows.

\textbf{CAMERA:} Slow lateral movement, low to the ground, following the stones before rising to Tlacaelel.

\textbf{TLACAELEL} (28) sits cross-legged, cloak folded neatly beside him. He marks figures with a charcoal stylus, pauses, erases, marks again. His movements are economical, habitual.

\textbf{CITLALI} (26) enters from the courtyard, carrying a painted clay cup of pulque. She sets it down within his reach but does not interrupt him.

She watches.

A beat.

\textbf{CITLALI}  
You count again.

Tlacaelel does not look up.

\textbf{TLACAELEL}  
Always.

\textbf{CITLALI}  
It is late.

\textbf{TLACAELEL}  
The numbers do not know that.

He finally looks up. His eyes are sharp but tired.

\textbf{CITLALI}  
What do they say tonight?

Tlacaelel gestures to the codex, then to the stones.

\textbf{TLACAELEL}  
That we take more men than the maize fields replace.  
That victories are getting heavier.  
That every campaign costs us more elders than boys.

Citlali sits opposite him, folding her legs beneath her.

\textbf{CITLALI}  
Your brother spoke at the market today.

Tlacaelel stiffens slightly.

\textbf{TLACAELEL}  
What did he say?

\textbf{CITLALI}  
That you measure the gods as if they were granaries.

Tlacaelel exhales slowly.

\textbf{TLACAELEL}  
I measure us.  
The gods can speak for themselves.

\textbf{CITLALI}  
They already have. Through the priests.

She studies him carefully.

\textbf{CITLALI}  
They say you doubt necessity.

\textbf{TLACAELEL}  
No. I doubt waste.

A distant drumbeat grows slightly louder, then fades.

\textbf{CITLALI}  
That difference will not matter to them.

Tlacaelel closes the codex halfway, thinking.

\textbf{TLACAELEL}  
It will matter when there are no men left to argue.

A knock at the door. Three sharp raps. Formal.

Citlali rises. Tlacaelel is already standing.

The door opens. \textbf{NEZAHUALCOYOTL} (30) enters, travel-dust on his sandals, wearing the blue-and-white of Texcoco. His expression is alert, amused, and wary.

He and Tlacaelel clasp forearms.

\textbf{NEZAHUALCOYOTL}  
Your uncle wants you at the council. Now.

\textbf{TLACAELEL}  
At this hour?

\textbf{NEZAHUALCOYOTL}  
Your arithmetic has learned how to walk.  
It reached the priests before I did.

Citlali inclines her head politely.

\textbf{CITLALI}  
You come as messenger or witness?

\textbf{NEZAHUALCOYOTL}  
Both, if I’m careful.

(to Tlacaelel)  
They say you want to make war seasonal.

\textbf{TLACAELEL}  
I want it survivable.

Nezahualcoyotl glances at the codex, at the stones.

\textbf{NEZAHUALCOYOTL}  
You’re counting what no one wants counted.

\textbf{TLACAELEL}  
Then someone must.

Citlali places a hand briefly on Tlacaelel’s arm.

\textbf{CITLALI}  
If you do this, do it completely.  
Half-measures invite knives.

Tlacaelel meets her eyes.

\textbf{TLACAELEL}  
I know.

He rolls the codex closed with care, gathers the counting stones into a pouch.

\textbf{TLACAELEL}  
War should not be a festival of endings.  
It should return.  
Like breath.

Nezahualcoyotl nods slowly.

\textbf{NEZAHUALCOYOTL}  
Then let’s go explain that to men who profit from endings.

Tlacaelel turns back once, looking at the room, the quiet order of it.

\textbf{TLACAELEL}  
(to Citlali)  
If I am not back before dawn—

\textbf{CITLALI}  
—then I will assume you are making trouble  
that requires daylight.

A faint smile passes between them.

Tlacaelel and Nezahualcoyotl exit into the night.

\textbf{CAMERA:} Lingers on Citlali alone. She kneels, opens the codex again, studies the figures Tlacaelel left unfinished.

The drums continue, distant and unresolved.

\textbf{CUT TO BLACK.}

\section*{Scene I-B}

\textbf{EXT. TOLTEC RUINS OUTSIDE THE CITY — DAWN}

Mist clings low to the ground. The ruins emerge slowly as the light strengthens: massive stone platforms fractured by roots, toppled columns half-swallowed by vines, colossal heads with eroded features staring past the present.

\textbf{CAMERA:} Wide establishing shot, then a slow push forward through drifting fog, revealing carved \textbf{GLYPHS} along a fallen wall. The figures show warriors gripping arms, hooking legs, forcing kneeling submission. No throats are cut. No bodies lie prone.

\textbf{TLACAELEL} walks alone among the stones. He wears no insignia. His sandals scrape softly against ancient steps worn smooth by centuries of feet.

He stops before a carved relief, brushing moss away with his fingers.

\textbf{CAMERA:} Close on his hand tracing a glyph of two figures locked in controlled struggle.

Footsteps. Fast. Bare.

\textbf{XOCHITL} (16) appears at the edge of the frame, having run in without ceremony. She is breathing hard but controlled, posture straight despite exertion. A small \textbf{BONE WHISTLE} hangs at her throat, bouncing lightly.

She waits for acknowledgment.

\textbf{XOCHITL}  
(formal)  
Tlacaelel. The emperor summons you.

Tlacaelel does not turn.

\textbf{TLACAELEL}  
What do you see here?

Xochitl blinks, caught off-guard.

\textbf{XOCHITL}  
These ruins?

\textbf{TLACAELEL}  
Yes.

She hesitates, scanning the stones.

\textbf{XOCHITL}  
The work of the Toltecs.  
Old temples. Old wars.

Tlacaelel shakes his head slightly.

\textbf{TLACAELEL}  
Old agreements.

He turns now, gestures toward the carvings.

\textbf{TLACAELEL}  
They fought often.  
But they expected to meet again.

Xochitl steps closer, studying the glyphs more carefully.

\textbf{XOCHITL}  
They are holding each other.

\textbf{TLACAELEL}  
They are ending motion, not life.

A pause. Morning birds begin to stir.

Tlacaelel notices the whistle at her neck.

\textbf{TLACAELEL}  
That—may I?

Xochitl lifts the cord over her head and places the whistle in his palm. It is skull-shaped, its mouth open in a permanent scream.

\textbf{TLACAELEL}  
Do you know its name?

\textbf{XOCHITL}  
Ehecachichtli.  
Death whistle.

\textbf{TLACAELEL}  
And what is it for?

\textbf{XOCHITL}  
To frighten enemies.  
To announce the presence of death.

Tlacaelel raises it to his lips.

\textbf{CAMERA:} Tight on his face. A breath.

He blows.

The sound is \textbf{PIERCING}, animal, inhuman. It echoes through the ruins, bouncing off stone. Birds explode from the trees. Somewhere deeper in the jungle, a monkey answers in alarm.

Xochitl flinches despite herself.

Tlacaelel lowers the whistle.

\textbf{TLACAELEL}  
Fear travels faster than any runner.

He returns it to her carefully.

\textbf{TLACAELEL}  
But fear that never arrives becomes memory.  
Memory lasts longer.

Xochitl reties the cord, thoughtful now.

\textbf{TLACAELEL}  
How many messages have you carried this moon?

\textbf{XOCHITL}  
Forty-three to Texcoco.  
Twenty-eight to Tlacopan.  
Eleven to Chalco.

\textbf{TLACAELEL}  
And how many returned unanswered?

She thinks.

\textbf{XOCHITL}  
Two.

\textbf{TLACAELEL}  
Good.

\textbf{XOCHITL}  
Good?

\textbf{TLACAELEL}  
It means most people still believe words matter.

He begins walking toward the path back to the city. Xochitl falls into step beside him without being told.

\textbf{TLACAELEL}  
Soon you will carry messages that sound wrong.  
That tell men when to fight, where to stop,  
and how to lose without disappearing.

\textbf{XOCHITL}  
If they refuse?

Tlacaelel does not answer immediately.

\textbf{TLACAELEL}  
Then the refusal becomes the message.

They reach the edge of the ruins. The city is visible in the distance now, sunlight catching on water and stone.

\textbf{TLACAELEL}  
Keep the whistle close.  
We will need a sound that means  
something new is beginning.

Xochitl nods, uncertain but attentive.

\textbf{CAMERA:} Wide shot as they walk away from the ruins together, the ancient stones receding into mist behind them.

The jungle hum rises with the day.

\textbf{CUT TO:}

\section*{Scene I}

\textbf{EXT. CAUSEWAY INTO TENOCHTITLAN — MORNING}

The causeway stretches across Lake Texcoco, a narrow spine of stone suspended over water alive with canoes. Morning light fractures across the surface of the lake, dazzling and indifferent.

\textbf{CAMERA:} A high, drifting shot follows the returning column of warriors from a distance, then descends gradually into their midst.

The warriors do not march. They arrive in uneven clusters, some supporting wounded companions, others walking alone. Shields are cracked. Feathered crests are missing. There is no chant, no rhythm imposed on their return.

Tlacaelel moves among them, not leading, not observing from afar. He walks at their pace.

A \textbf{SCRIBE}, TLAMACAZQUI (50s), keeps several steps behind him, carrying thin wooden tablets and a charcoal stylus. He records quietly, without comment.

They pass a group of women waiting at the causeway’s edge. One recognizes a face in the crowd and steps forward, relief collapsing into restraint. Another scans the returning men and then turns away, already understanding.

Tlacaelel stops near a \textbf{YOUNG WARRIOR} (18), seated on the stone edge with his feet in the water. He scrubs dried blood from his hands. His \textbf{MACUAHUITL} lies beside him, its wooden body intact but missing two obsidian blades.

\textbf{TLACAELEL}  
Your first campaign?

The young warrior does not look up.

\textbf{YOUNG WARRIOR}  
First that lasted longer than a morning.

Tlacaelel crouches, eye level now.

\textbf{TLACAELEL}  
You’re injured?

\textbf{YOUNG WARRIOR}  
No.  
Ixcatzin is. And Tezcatl—

He stops himself.

\textbf{YOUNG WARRIOR}  
—Tezcatl won’t be crossing back.

Tlacaelel glances at the water, at the ripples erasing reflections.

\textbf{TLACAELEL}  
How many did you take?

The young warrior finally looks up, confused.

\textbf{YOUNG WARRIOR}  
Take?

\textbf{TLACAELEL}  
Captives.

The young warrior shakes his head.

\textbf{YOUNG WARRIOR}  
We killed twelve.  
They didn’t break fast enough.

Tlacaelel nods once. Not approval. Acknowledgment.

\textbf{TLACAELEL}  
And how does that feel?

The young warrior laughs once, short and humorless.

\textbf{YOUNG WARRIOR}  
Empty.  
Like we spent something we can’t earn again.

Tlacaelel absorbs this without reaction.

\textbf{TLACAELEL}  
If I told you there was another way to fight—  
one that leaves fewer men behind in the mud—  
would you listen?

The young warrior studies him now, wary.

\textbf{YOUNG WARRIOR}  
Would it keep Ixcatzin alive?

Tlacaelel does not answer immediately.

\textbf{TLACAELEL}  
It would make his death rarer.

The young warrior nods slowly. That is enough.

Tlacaelel rises and continues walking.

TLAMACAZQUI catches up beside him.

\textbf{TLAMACAZQUI}  
(quiet)  
They will repeat that conversation.  
Word spreads faster than orders.

\textbf{TLACAELEL}  
Good.

\textbf{TLAMACAZQUI}  
The priests will hear it.

\textbf{TLACAELEL}  
Then they can count with us.

They approach the city gate. The towers loom, heavy with authority. Drums sound again, closer now, but still restrained.

\textbf{CAMERA:} Lingers on the young warrior as he watches Tlacaelel disappear into the crowd. He looks down at his broken weapon, then out across the lake.

He does not throw the macuahuitl away.

\textbf{CUT TO:}

\section*{Scene II-A}

\textbf{INT. EMPEROR’S PRIVATE GARDEN — LATE MORNING}

The garden is terraced and enclosed, hidden from the city’s noise by high stone walls softened with climbing flowers. Water runs through narrow channels carved into the stone, feeding pools of lilies and reeds. Everything here is cultivated, controlled, alive by design.

\textbf{CAMERA:} Begins high, looking down on the geometry of the garden, then descends slowly until individual leaves and droplets of water become visible.

\textbf{ITZCOATL} (40s), the emperor, stands with his hands clasped behind his back, studying a bed of blooming \textbf{XOCHITL}—flowers carefully spaced, pruned, thriving. He does not turn as Tlacaelel is escorted in, then dismissed.

Tlacaelel stops several paces behind him. He waits.

A long silence. Water continues to flow.

\textbf{ITZCOATL}  
Do you know why I like this garden?

Tlacaelel answers carefully.

\textbf{TLACAELEL}  
Because nothing here grows by accident.

Itzcoatl inclines his head, still facing the flowers.

\textbf{ITZCOATL}  
Because everything here would choke itself to death  
if I let it.

He finally turns.

\textbf{ITZCOATL}  
They tell me you want to change war.

\textbf{TLACAELEL}  
I want to change what it leaves behind.

\textbf{ITZCOATL}  
The priests say you are questioning necessity.

\textbf{TLACAELEL}  
No. I am questioning appetite.

Itzcoatl studies him with something like amusement.

\textbf{ITZCOATL}  
Careful. Hunger is the one thing emperors  
are allowed to understand.

Footsteps approach from the far terrace.

\textbf{MOCTEZUMA} (30), armored, dust still clinging to his sandals, enters without ceremony. He bows briefly to Itzcoatl, then looks directly at Tlacaelel.

\textbf{MOCTEZUMA}  
The warriors say you criticize their courage.

\textbf{TLACAELEL}  
I criticize their attrition.

Moctezuma’s jaw tightens.

\textbf{MOCTEZUMA}  
Men are not maize. You cannot replant them.

\textbf{TLACAELEL}  
Exactly.

Silence. Water trickles. A bird calls.

\textbf{ITZCOATL}  
Moctezuma thinks war ends when the enemy kneels.

\textbf{TLACAELEL}  
And I think war ends when the enemy vanishes.

\textbf{MOCTEZUMA}  
That is victory.

\textbf{TLACAELEL}  
That is amnesia.

Moctezuma steps forward.

\textbf{MOCTEZUMA}  
You would teach our enemies how to survive us.

\textbf{TLACAELEL}  
I would teach them how to need us.

Itzcoatl raises a hand. Moctezuma stops.

\textbf{ITZCOATL}  
Enough.

He plucks a single flower from the bed—a red dahlia—and turns it in his fingers.

\textbf{ITZCOATL}  
Xochitl.  
Beautiful because it is tended.  
Dead if neglected.  
Dangerous if allowed to spread.

He holds it out to Tlacaelel.

\textbf{ITZCOATL}  
You want to make war like this.

\textbf{TLACAELEL}  
Seasonal.  
Bounded.  
Expected.

Itzcoatl nods slowly.

\textbf{ITZCOATL}  
And if it fails?

\textbf{TLACAELEL}  
Then it fails slowly enough for us to adapt.

Itzcoatl looks to Moctezuma.

\textbf{ITZCOATL}  
Leave us.

Moctezuma hesitates, then bows and exits, eyes never leaving Tlacaelel.

A heavier silence settles.

\textbf{ITZCOATL}  
He will rule after me.

\textbf{TLACAELEL}  
I know.

\textbf{ITZCOATL}  
And he will remember this conversation  
no matter how it ends.

\textbf{TLACAELEL}  
So will the city.

Itzcoatl considers this, then places the flower in Tlacaelel’s hand.

\textbf{ITZCOATL}  
Go to the council.  
Propose your war of flowers.  
I will support you—once.

\textbf{TLACAELEL}  
That is all I ask.

\textbf{ITZCOATL}  
Good.  
Because men who change structures  
rarely get a second invitation.

Tlacaelel bows deeply.

\textbf{CAMERA:} Holds on Itzcoatl alone as Tlacaelel exits, then slowly returns to the flower bed, now missing one bloom.

Water continues to flow.

\textbf{CUT TO:}

\section*{Scene II}

\textbf{INT. COUNCIL CHAMBER — LATE MORNING}

The chamber is broad and low, its ceiling supported by thick wooden beams darkened by generations of smoke. Light enters from a narrow opening high above, falling in a single pale shaft that moves slowly across the floor as the sun advances.

Stone benches line the walls. Elders, nobles, war captains, and priests sit in deliberate spacing. No one rushes to speak. Authority here is measured by patience.

\textbf{CAMERA:} Static at first, holding the geometry of the room, then slowly tracking the movement of the light across faces as voices emerge.

At the far end, the \textbf{EMPEROR} sits upright, listening. To one side stands \textbf{TIZOC}, High Priest of Huitzilopochtli, his expression composed and watchful. \textbf{MOCTEZUMA} sits among the captains, arms folded, unreadable.

Tlacaelel stands near the center, unadorned, hands empty.

A priest finishes speaking.

\textbf{PRIEST}  
The signs were favorable.  
The offerings were accepted.  
The sun rose as promised.

A low murmur of assent moves through the chamber, brief and controlled.

Tlacaelel waits for it to pass.

\textbf{TLACAELEL}  
The sun rose because the world was not emptied yesterday.

The chamber stills.

\textbf{TIZOC}  
Choose your words carefully.

\textbf{TLACAELEL}  
I am.

He steps forward half a pace, enough to claim attention without presumption.

\textbf{TLACAELEL}  
We fight well.  
Too well.

A few captains shift, offended but curious.

\textbf{TLACAELEL}  
Our enemies die where they stand.  
Their sons do not return.  
Their fields go untilled.  
Their names disappear.

An elder leans forward.

\textbf{ELDER}  
You speak as if this were loss.

\textbf{TLACAELEL}  
It is waste.

The word lands sharply. Administrative. Unemotional.

\textbf{MOCTEZUMA}  
Waste is the price of victory.

Tlacaelel turns to him.

\textbf{TLACAELEL}  
Victory that consumes its future  
is only delay.

Tizoc rises slowly.

\textbf{TIZOC}  
The gods require nourishment.

\textbf{TLACAELEL}  
Yes.  
But they require continuity more.

A murmur. Unease now.

\textbf{TLACAELEL}  
A dead enemy feeds the sun once.  
A living enemy feeds it again and again—  
through fear, obligation, and return.

\textbf{ELDER}  
You would spare them?

\textbf{TLACAELEL}  
I would keep them.

Moctezuma exhales, incredulous.

\textbf{MOCTEZUMA}  
You propose war without endings.

Tlacaelel meets his gaze.

\textbf{TLACAELEL}  
I propose war with rules.

Silence. The light shifts slightly across the floor.

\textbf{TLACAELEL}  
Battles fought at agreed fields.  
At agreed times.  
With warriors named in advance.

\textbf{TLACAELEL}  
Victory counted in captives, not bodies.  
Those taken will live.  
Those who fall will be few.

A councilor laughs quietly, without humor.

\textbf{COUNCILOR}  
You would turn war into ceremony.

\textbf{TLACAELEL}  
I would turn it into memory.

The Emperor finally speaks.

\textbf{EMPEROR}  
And if they refuse these terms?

Tlacaelel does not hesitate.

\textbf{TLACAELEL}  
Then we return to the old way.  
And one of us disappears.

The bluntness settles the room.

\textbf{TIZOC}  
You risk teaching the people  
that the gods can be managed.

\textbf{TLACAELEL}  
We have always managed them.  
We simply pretended otherwise.

A long pause.

The Emperor studies the faces around him, then Tlacaelel.

\textbf{EMPEROR}  
This council will consider your proposal.

Tlacaelel bows, neither relieved nor triumphant.

\textbf{CAMERA:} Holds on Moctezuma’s face as the meeting dissolves—calculating, unsettled.

The shaft of sunlight reaches the far wall.

\textbf{CUT TO:}

\section*{Scene III}

\textbf{EXT. WARRIORS’ TRAINING YARD — AFTERNOON}

A wide expanse of packed earth bordered by low stone walls. Training posts stand at measured intervals. Weapons racks line one side: shields, macuahuitl, spears. The air carries dust, sweat, and the dull rhythm of repeated impact.

\textbf{CAMERA:} Handheld but restrained, close to the bodies—shoulders, feet, breath—never drifting into spectacle.

A mixed group of \textbf{YOUNG WARRIORS}, including several women, trains under the supervision of \textbf{VETERAN INSTRUCTORS}. Movements are precise, repetitive, exhausting.

Tlacaelel stands at the edge with TLAMACAZQUI, observing without interruption.

Two trainees face each other: \textbf{YACANEX} (19, male) and \textbf{MALINAL} (20, female). Both wear light padding, their macuahuitl fitted with blunted obsidian edges near the hilt.

\textbf{INSTRUCTOR}  
Again.  
Control first. Force second.

They engage. Yacanex swings too wide. Malinal steps inside his guard, hooks his leg, and forces him down onto one knee. Her blade stops at his throat.

She freezes, waiting.

\textbf{INSTRUCTOR}  
Hold it.

He steps between them.

\textbf{INSTRUCTOR}  
That is the moment.  
Not the cut.  
The decision.

Yacanex breathes hard, frustrated.

\textbf{YACANEX}  
In battle, he would already be dead.

\textbf{INSTRUCTOR}  
In battle, you would already be dead  
if you hesitated that long.

The instructor looks to Malinal.

\textbf{INSTRUCTOR}  
You felt it.

\textbf{MALINAL}  
Yes.

\textbf{INSTRUCTOR}  
Say it.

\textbf{MALINAL}  
I could end him.

\textbf{INSTRUCTOR}  
And you did not.

The instructor steps back.

\textbf{INSTRUCTOR}  
Again.

They reset.

Nearby, another pair trains—one male, one female—grappling without weapons, working leverage and balance rather than strikes.

Tlacaelel watches closely.

\textbf{TLAMACAZQUI}  
They are confused.

\textbf{TLACAELEL}  
They are thinking.

A veteran captain approaches, helmet under his arm.

\textbf{CAPTAIN}  
They will hesitate in real combat.

\textbf{TLACAELEL}  
They already hesitate.  
The difference is whether hesitation has shape.

The captain watches Malinal force Yacanex off balance again, cleaner this time.

\textbf{CAPTAIN}  
Women in the yard will unsettle the old houses.

\textbf{TLACAELEL}  
So will losing sons.

The captain says nothing.

A sharp sound cuts through the yard.

\textbf{CAMERA:} Whip-pan to a rack where a modified macuahuitl has snapped—obsidian shards scattered in the dirt.

An apprentice stares at it, shaken.

\textbf{APPRENTICE}  
It broke.

\textbf{INSTRUCTOR}  
Pick it up.

The apprentice kneels, gathering the shards carefully, as if they might cut him even now.

Tlacaelel steps forward.

\textbf{TLACAELEL}  
(to the yard)  
Stop.

Training halts. All eyes turn to him.

\textbf{TLACAELEL}  
What you are learning is not how to win quickly.  
It is how to win without erasing the field.

Murmurs ripple—confusion, curiosity, resistance.

\textbf{TLACAELEL}  
If you want to kill, do it cleanly and be done.  
If you want to fight again, learn restraint.

A young woman at the back speaks up.

\textbf{YOUNG WARRIOR (FEMALE)}  
And if the enemy doesn’t restrain themselves?

Tlacaelel meets her gaze.

\textbf{TLACAELEL}  
Then you will be better trained  
when the rules break.

A beat.

\textbf{TLACAELEL}  
Resume.

Training restarts, slower now, more deliberate.

\textbf{CAMERA:} Pulls back as bodies move in controlled rhythm, the dust rising evenly rather than chaotically.

Tlacaelel watches, expression unreadable.

\textbf{CUT TO:}

\section*{Scene III-A}

\textbf{INT. PRIESTS’ TRAINING CHAMBER — LATE AFTERNOON}

The chamber is cool and dim, carved stone absorbing the heat of the day. Copal incense hangs thick in the air, its sweetness edged with bitterness. Painted walls depict the \textbf{FIVE SUNS}: worlds born, destroyed, reborn. Each ending is violent. Each rebirth conditional.

\textbf{CAMERA:} Slow pan across the murals, lingering on the moment of collapse in each cycle, before settling on the present Sun—unfinished, still in motion.

A group of \textbf{YOUNG PRIESTS} kneels in ordered rows, chanting softly as they trace calendar glyphs onto bark-paper sheets.

At the front stands \textbf{TIZOC} (50s), High Priest, composed and severe. His presence fills the room without effort.

\textbf{TIZOC}  
The sun does not move because it is kind.  
It moves because it is fed.

He gestures to a glyph showing blood rising as smoke.

\textbf{TIZOC}  
Tonalli ascends.  
Breath becomes light.  
This is not metaphor.

The young priests nod, repeat the chant.

Tlacaelel enters quietly, stopping at the threshold. He waits.

Tizoc continues teaching without acknowledging him.

\textbf{TIZOC}  
Four suns failed because people mistook habit for stability.  
They assumed the world would forgive delay.

A young priest looks up.

\textbf{YOUNG PRIEST}  
What happens if the offerings are late?

Tizoc’s gaze sharpens.

\textbf{TIZOC}  
Then the sun remembers its hunger.

He finally turns toward Tlacaelel.

\textbf{TIZOC}  
You have come to adjust the calendar.

\textbf{TLACAELEL}  
I have come to preserve it.

Tizoc dismisses the young priests with a flick of his hand. They bow and exit in silence.

The two men stand facing one another beneath the painted suns.

\textbf{TIZOC}  
You want to ration blood.

\textbf{TLACAELEL}  
I want to prevent famine.

\textbf{TIZOC}  
The gods do not starve the way men do.

\textbf{TLACAELEL}  
Men do the starving for them.

A charged silence.

\textbf{TIZOC}  
You are dangerous.

\textbf{TLACAELEL}  
So is arithmetic.

Tizoc steps closer, lowering his voice.

\textbf{TIZOC}  
If you teach the people that sacrifice can be planned,  
they will learn to delay it.

\textbf{TLACAELEL}  
If we do not plan it,  
we will run out of people to teach.

Tizoc studies him carefully now, reassessing.

\textbf{TIZOC}  
You are not questioning the gods.

\textbf{TLACAELEL}  
No.

\textbf{TIZOC}  
You are questioning us.

\textbf{TLACAELEL}  
Yes.

Tizoc turns back to the mural of the Fifth Sun.

\textbf{TIZOC}  
This world ends by earthquake, according to the count.  
Not today.  
But it will end.

\textbf{TLACAELEL}  
Every system ends.

\textbf{TIZOC}  
And yet you build one that depends on restraint.

\textbf{TLACAELEL}  
Because unrestrained systems end faster.

Tizoc considers this, then allows himself a thin smile.

\textbf{TIZOC}  
You are either the most necessary man in this city  
or the one who will bring it down.

\textbf{TLACAELEL}  
Those roles often overlap.

A beat.

\textbf{TIZOC}  
Proceed.  
But understand this: if the sun falters,  
your name will be spoken with knives.

\textbf{TLACAELEL}  
If the sun falters,  
names will be the least of our problems.

They hold each other’s gaze.

\textbf{CAMERA:} Tilts upward, framing both men beneath the painted Sun, its incomplete lines suggesting a future not yet decided.

\textbf{CUT TO:}

\section*{Scene IV}

\textbf{EXT. JUNGLE RELAY STATION — DAY}

A small stone structure sits at the edge of a dense jungle path. Moss creeps along its walls. A wooden marker carved with empire glyphs stands nearby, weathered but maintained.

\textbf{CAMERA:} Low tracking shot following bare feet striking stone, mud, roots—fast, light, practiced.

\textbf{XOCHITL} bursts into frame, breath steady despite the distance she has covered. Sweat darkens her tunic. She slows only when she reaches the shade of the station.

\textbf{YAOTL} (40s), the station keeper, waits with a gourd of water already extended. His left leg bears the mark of an old injury; he favors it without complaint.

Xochitl drinks deeply, hands him a wrapped bark-paper message.

\textbf{XOCHITL}  
From Tlacaelel.  
For Tlaxcala.

Yaotl breaks the seal, scans the glyphs.

\textbf{YAOTL}  
“Flowering.”  
That is the word he chose?

\textbf{XOCHITL}  
He says words shape how people listen.

Yaotl snorts softly.

\textbf{YAOTL}  
I was at Azcapotzalco.  
No one listened. They survived or they didn’t.

Xochitl tightens the strap of her message pouch.

\textbf{XOCHITL}  
And now?

Yaotl gestures toward the jungle.

\textbf{YAOTL}  
Now we run until someone answers.

Movement in the trees. Another \textbf{RUNNER} emerges—Tlaxcalan markings on his sash. He is younger, wary, measuring the space before stepping closer.

\textbf{TLAXCALAN RUNNER}  
Formal message.  
For Tlacaelel of Tenochtitlan.

He hands over a sealed packet.

Yaotl opens it, reads, then looks up at Xochitl.

\textbf{YAOTL}  
They are listening.

\textbf{XOCHITL}  
Listening or preparing?

\textbf{YAOTL}  
In my experience, the difference is small.

He hands her the reply.

\textbf{YAOTL}  
They want assurances.  
Boundaries.  
Names of fields.  
Days counted in advance.

Xochitl exhales, half relief, half dread.

\textbf{XOCHITL}  
Then it’s real.

\textbf{YAOTL}  
It was real when the first man decided  
not to kill someone who expected it.

The Tlaxcalan runner drinks from the gourd, eyes on Xochitl.

\textbf{TLAXCALAN RUNNER}  
If this is a trick—

\textbf{XOCHITL}  
—then it’s a slow one.  
And slow tricks give people time to notice.

She turns back toward the path to Tenochtitlan.

\textbf{YAOTL}  
You’re not done running today.

\textbf{XOCHITL}  
I know.

She ties her sandals tighter.

\textbf{XOCHITL}  
The gods made me fast for a reason.

\textbf{YAOTL}  
The gods made you young.

Xochitl smiles briefly.

\textbf{XOCHITL}  
Then let’s use that before they change their minds.

She disappears into the jungle, movement swallowed by green.

\textbf{CAMERA:} Holds on the relay station as jungle sounds reclaim the space. The Tlaxcalan runner watches the path she took, thoughtful.

\textbf{CUT TO:}

\section*{Scene V}

\textbf{INT. MOCTEZUMA’S QUARTERS — NIGHT}

The room is long and severe. Stone walls are hung with captured shields, enemy banners, and weapons taken in battle. Each object is placed deliberately, not cluttered, not ornamental. This is a space meant to remember force.

\textbf{CAMERA:} Slow pan across the trophies—splintered shields, cracked macuahuitl, jaguar pelts—before settling on \textbf{MOCTEZUMA} (30), seated on a woven mat, removing dried blood from his forearms with a cloth.

Around him sit three \textbf{VETERAN CAPTAINS}, scarred, disciplined. Near Moctezuma’s right hand stands \textbf{CUAUHTÉMOC} (18), alert, ambitious, absorbing every word.

A silence hangs before anyone speaks.

\textbf{CAPTAIN ONE}  
The men are unsettled.

Moctezuma does not look up.

\textbf{MOCTEZUMA}  
Good.  
Unsettled men listen.

\textbf{CAPTAIN TWO}  
They are being told not to finish fights.

Moctezuma pauses his cleaning.

\textbf{MOCTEZUMA}  
They are being told to hesitate.

\textbf{CAPTAIN THREE}  
Hesitation gets warriors killed.

Moctezuma looks up now.

\textbf{MOCTEZUMA}  
So does exhaustion.  
So does arrogance.

Cuauhtémoc steps forward.

\textbf{CUAUHTÉMOC}  
Let me lead a raid.  
Just one.  
Show them the old way still works.

Moctezuma rises slowly, towering over him.

\textbf{MOCTEZUMA}  
No.

Cuauhtémoc stiffens.

\textbf{CUAUHTÉMOC}  
Why not?

\textbf{MOCTEZUMA}  
Because that is exactly what Tlacaelel wants.  
Someone to break the pattern  
so he can say the pattern was necessary.

The captains exchange looks.

\textbf{CAPTAIN ONE}  
Then what do we do?

Moctezuma walks to the wall, touches a shield marked with Tlaxcalan glyphs.

\textbf{MOCTEZUMA}  
We watch.

\textbf{CAPTAIN TWO}  
Watch what?

\textbf{MOCTEZUMA}  
A garden.

They turn to him.

\textbf{MOCTEZUMA}  
Gardens look strong when they bloom.  
Soft, even.

He taps the shield.

\textbf{MOCTEZUMA}  
But the first frost kills them all at once.

Cuauhtémoc frowns.

\textbf{CUAUHTÉMOC}  
What if it works?

Moctezuma turns, sharp.

\textbf{MOCTEZUMA}  
Then we will have taught our enemies  
how to survive us.

A pause.

\textbf{CUAUHTÉMOC}  
Is that so terrible?

Moctezuma studies him—measuring, weighing.

\textbf{MOCTEZUMA}  
Empires do not fall when they are defeated.  
They fall when they become unnecessary.

He steps closer to Cuauhtémoc.

\textbf{MOCTEZUMA}  
Remember that.

The captains bow their heads.

\textbf{MOCTEZUMA}  
Leave me.

They exit, one by one. Cuauhtémoc hesitates, then follows.

Moctezuma remains alone.

\textbf{CAMERA:} Moves in close as he places the cleaned cloth beside him, hands now still.

He looks at the trophies again—not with pride, but calculation.

\textbf{MOCTEZUMA}  
(quietly, to himself)  
War is hunger.

He extinguishes a torch.

Darkness takes the room.

\textbf{CUT TO:}

\section*{Scene VI}

\textbf{INT. OBSIDIAN WORKSHOP — DAY}

Light filters in through a high opening, catching suspended dust and turning it into a slow, glittering haze. The air is sharp with stone and sweat. Every surface bears the marks of work: grooves cut into wood, black flakes embedded in the earth floor.

\textbf{CAMERA:} Close, intimate—hands, tools, edges—never pulling wide enough to romanticize the space.

At a low bench sits \textbf{IXTLIL} (60s), master weapon-crafter, his left eye clouded from an old injury. His hands move with unhurried precision, striking obsidian at exact angles. Each tap produces a clean, ringing click.

Several \textbf{APPRENTICES} work nearby, quieter, less certain.

Tlacaelel enters and waits. He does not interrupt the rhythm.

A blade fractures cleanly. Ixtlil inspects it, nods, then finally looks up.

\textbf{IXTLIL}  
You’re late.

\textbf{TLACAELEL}  
I didn’t want to break the pattern.

Ixtlil snorts.

\textbf{IXTLIL}  
Patterns break whether you want them to or not.

He gestures toward a rack of finished \textbf{MACUAHUITL}. Some look familiar. Others have visibly altered profiles—edges sharpened only toward the tip, inner sections dulled.

\textbf{IXTLIL}  
These are your ideas.

Tlacaelel steps closer, lifts one carefully.

\textbf{TLACAELEL}  
They’re your work.

\textbf{IXTLIL}  
They’re compromises.

He takes the weapon back, runs a thumb just shy of the edge.

\textbf{IXTLIL}  
Obsidian doesn’t forgive intention.  
It cuts or it shatters.

\textbf{TLACAELEL}  
So do men.

Ixtlil studies him.

\textbf{IXTLIL}  
You want weapons that wound without finishing.  
That frighten without ending.  
That requires discipline.

\textbf{TLACAELEL}  
Yes.

\textbf{IXTLIL}  
Then you’re making war harder.

Tlacaelel nods.

\textbf{TLACAELEL}  
I am.

An apprentice approaches, holding a broken blade.

\textbf{APPRENTICE}  
This failed in training.

Ixtlil takes it, examines the fracture.

\textbf{IXTLIL}  
It didn’t fail.  
It was treated like a toy.

He looks at Tlacaelel.

\textbf{IXTLIL}  
Warriors respect weapons that punish mistakes.  
You’re asking them to think instead.

\textbf{TLACAELEL}  
Thinking is cheaper than funerals.

Ixtlil allows himself a thin smile.

\textbf{IXTLIL}  
My grandfather made blades for the Toltecs.  
He told me they had weapons for ritual combat—  
edges shaped for capture.

\textbf{TLACAELEL}  
What happened to them?

\textbf{IXTLIL}  
The Toltecs ended.  
And people decided the blades were myths.

Tlacaelel sets the broken obsidian carefully on the bench.

\textbf{TLACAELEL}  
Or warnings.

Ixtlil watches him for a long moment, then hands a finished weapon to an apprentice.

\textbf{IXTLIL}  
Make ten more like this.  
Slowly.

He turns back to Tlacaelel.

\textbf{IXTLIL}  
If this works, it won’t be because the weapons are clever.  
It will be because the men holding them are.

\textbf{TLACAELEL}  
That’s the hope.

\textbf{IXTLIL}  
Hope is brittle.

\textbf{TLACAELEL}  
So is obsidian.  
And yet it still cuts the world open.

A beat.

\textbf{CAMERA:} Holds on the craftsmen resuming their rhythm as Tlacaelel exits, the sound of stone striking stone steady, methodical, unforgiving.

\textbf{CUT TO:}

\section*{Scene VII}

\textbf{INT. TLACAELEL’S HOME — NIGHT}

The central room has been cleared and prepared for guests. Woven mats form a square around low serving platters. Oil lamps flicker, throwing warm light over painted bowls and steam rising from food.

\textbf{CAMERA:} Slow, deliberate movement, tracking the placement of each dish as if inventory were being taken: tamales wrapped in corn husks, bowls of thick mole dark with cacao and chile, roasted turkey, baskets of maize bread. Painted gourds filled with frothed \textbf{XOCOATL} are set carefully within reach.

\textbf{CITLALI} moves quietly, directing younger attendants with subtle gestures. Her face is composed, hospitable, alert.

\textbf{XICOTENCATL} (45), Tlaxcalan war chief, enters with two aides. His posture is guarded but respectful. He wears no armor, only a cloak marked with the glyphs of his city.

Tlacaelel rises to greet him. They incline their heads—neither bowing fully.

\textbf{TLACAELEL}  
You honor my house.

\textbf{XICOTENCATL}  
I enter it carefully.

They sit opposite one another. Nezahualcoyotl positions himself slightly to the side, mediator by stance alone.

Food is served. For a moment, no one speaks.

\textbf{CAMERA:} Close on hands breaking bread, steam fogging painted surfaces, the ritual of eating before words.

Xicotencatl tastes the chocolate, then nods once.

\textbf{XICOTENCATL}  
Bitter.  
You have not softened it for guests.

\textbf{CITLALI}  
We serve it as it is meant to be taken.

Xicotencatl meets her eyes, amused despite himself.

\textbf{XICOTENCATL}  
Good.

He turns back to Tlacaelel.

\textbf{XICOTENCATL}  
Your message says you want to fight us  
without trying to erase us.

\textbf{TLACAELEL}  
I want to fight you  
so that you remain.

A flicker of something like approval crosses Xicotencatl’s face, then vanishes.

\textbf{XICOTENCATL}  
In Tlaxcala, we teach our children  
never to trust a Mexica promise.

\textbf{TLACAELEL}  
And in Tenochtitlan, we teach ours  
never to underestimate Tlaxcala.

Nezahualcoyotl intervenes smoothly.

\textbf{NEZAHUALCOYOTL}  
This is not surrender.  
It is structure.

Xicotencatl leans back slightly.

\textbf{XICOTENCATL}  
You want agreed fields.  
Agreed days.  
Agreed limits.

\textbf{TLACAELEL}  
Yes.

\textbf{XICOTENCATL}  
And captives?

\textbf{TLACAELEL}  
Counted.  
Returned when the calendar allows.

Silence. The lamps crackle.

\textbf{XICOTENCATL}  
You are asking us to trust restraint  
from men who built an empire on hunger.

Tlacaelel nods.

\textbf{TLACAELEL}  
Yes.

\textbf{XICOTENCATL}  
Why should we?

Tlacaelel does not answer immediately. He gestures to the food.

\textbf{TLACAELEL}  
Because we are eating together  
instead of sharpening knives.

Xicotencatl considers this, then sets his cup down.

\textbf{XICOTENCATL}  
My council will say this is slow death.

\textbf{TLACAELEL}  
Then tell them the alternative  
is quick extinction.

Citlali watches Xicotencatl carefully.

\textbf{CITLALI}  
If we betray this,  
you will be justified in breaking it.

Xicotencatl studies her, then Tlacaelel.

\textbf{XICOTENCATL}  
And if it works?

Tlacaelel answers softly.

\textbf{TLACAELEL}  
Then both our children  
will learn how to fight without disappearing.

A long pause.

\textbf{XICOTENCATL}  
I will bring this to the Four Hundred.  
Do not mistake that for agreement.

\textbf{TLACAELEL}  
I won’t.

Xicotencatl rises. His aides follow.

At the doorway, he stops.

\textbf{XICOTENCATL}  
One thing more.

Tlacaelel looks up.

\textbf{XICOTENCATL}  
If this fails,  
you will have taught us how to endure you longer.

\textbf{TLACAELEL}  
And you will have taught us  
that endurance cuts both ways.

Xicotencatl nods once and leaves.

The room remains still.

Citlali exhales quietly.

\textbf{CITLALI}  
He will agree.

\textbf{TLACAELEL}  
I know.

\textbf{CITLALI}  
That troubles you.

Tlacaelel looks at the remaining food, untouched now.

\textbf{TLACAELEL}  
Because once restraint becomes habit,  
the first person to abandon it  
wins everything.

Nezahualcoyotl meets his gaze.

\textbf{NEZAHUALCOYOTL}  
Then your task is not to build peace.  
It is to make betrayal expensive.

Tlacaelel nods slowly.

\textbf{CUT TO:}

\section*{Scene VIII}

\textbf{EXT. DESIGNATED FLOWER WAR FIELD — MIDDAY}

A broad, open plain bordered by low hills. No villages nearby. No crops. Markers of stone and wood define the perimeter. Banners bearing Mexica and Tlaxcalan glyphs stand opposite one another, aligned with ritual precision.

\textbf{CAMERA:} Wide establishing shot, symmetrical, almost ceremonial. Then the symmetry breaks as bodies enter the frame.

Warriors assemble in ordered ranks. Faces are painted. Armor is lighter than in conquest battles. Weapons are present but restrained—modified macuahuitl, spears without barbs.

Drums begin. Not a charge rhythm. A count.

\textbf{TLATOANI} (30s), veteran Mexica warrior, stands with \textbf{YACANEX} nearby. Across the field, Tlaxcalan warriors mirror their formation.

\textbf{TLATOANI}  
(low)  
Remember the field.

\textbf{YACANEX}  
I see it.

\textbf{TLATOANI}  
Remember it when you don’t want to.

The drums stop.

A single \textbf{DEATH WHISTLE} sounds—short, sharp.

The lines advance.

\textbf{CAMERA:} Drops into the movement. Breath, footwork, controlled chaos. No sweeping heroics—only proximity and consequence.

Yacanex engages a Tlaxcalan warrior, \textbf{OPOCHTZIN} (20s), agile, alert. They circle, test distance.

Steel does not clash. Obsidian hums.

Opochtzin feints high. Yacanex reacts too quickly, overcommits. Opochtzin nearly takes his leg.

\textbf{TLATOANI}  
(low, urgent)  
Slow!

Yacanex regains footing. He remembers. He hooks instead of slashes. Opochtzin stumbles.

They grapple. Dirt. Breath in each other’s faces.

Yacanex gains leverage, presses the blade edge to Opochtzin’s throat—but stops.

Time stretches.

\textbf{OPOCHTZIN}  
(breathless)  
Neomiquiz.

Yacanex hesitates.

\textbf{TLATOANI}  
Correct him.

Yacanex swallows.

\textbf{YACANEX}  
Nicochi.

Opochtzin exhales, a strange half-laugh.

\textbf{OPOCHTZIN}  
New words.

Yacanex steps back, keeping his blade close but controlled.

\textbf{YACANEX}  
Can you walk?

Opochtzin tests his weight.

\textbf{OPOCHTZIN}  
Yes.

\textbf{YACANEX}  
Then walk.

They move together toward the collection point where other captives are being gathered—not bound, not struck, guarded with alert restraint.

Nearby, another engagement ends badly—a warrior strikes too hard. Blood sprays. A man collapses.

A beat of shock.

The drums falter, then resume, slower.

\textbf{CAMERA:} Finds Tlacaelel at the field’s edge, watching intently, jaw tight. Tizoc stands beside him, unreadable.

\textbf{TIZOC}  
One death.

\textbf{TLACAELEL}  
One.

\textbf{TIZOC}  
They will count it.

\textbf{TLACAELEL}  
They should.

The battle continues. Captures outnumber kills. The rhythm stabilizes.

Across the field, warriors shout names—not insults, not war cries, but identifications.

\textbf{TLAXCALAN WARRIOR (O.S.)}  
Opochtzin!

\textbf{OPOCHTZIN}  
(shouting back)  
Next count! Same field!

Yacanex looks at him, confused but steady.

\textbf{YACANEX}  
You know when we fight again?

\textbf{OPOCHTZIN}  
Everyone knows.  
That’s the point.

The final death whistle sounds—longer this time.

The fighting stops.

Bodies remain standing. The field holds.

\textbf{CAMERA:} Pulls back slowly, revealing both sides still present, wounded but extant, the space between them intact.

Tlacaelel closes his eyes briefly.

Not relief.

Calculation.

\textbf{CUT TO:}

\section*{Scene IX}

\textbf{INT. TENOCHTITLAN — VARIOUS — EVENING}

\textbf{CAMERA:} A sequence of short, intercut moments. The city absorbs news the way stone absorbs heat—slowly, unevenly.

\textbf{— MARKET}

Vendors pack away goods. A woman recounts the day to another, her voice low.

\textbf{MARKET WOMAN}  
They came back early.

\textbf{OTHER WOMAN}  
Early how?

\textbf{MARKET WOMAN}  
Standing.

They exchange a look that carries both relief and suspicion.

\textbf{— CANAL}

A canoe glides through narrow water. Two old men fish silently. One speaks without turning.

\textbf{FISHERMAN}  
My grandson fought today.

\textbf{OTHER FISHERMAN}  
Did he return?

A pause.

\textbf{FISHERMAN}  
Yes.

They continue fishing. No celebration. No prayer. Just continuation.

\textbf{— TEMPLE STEPS}

Young acolytes whisper among themselves. One glances upward toward the shrine.

\textbf{YOUNG ACOLYTE}  
There was only one offering.

\textbf{ANOTHER}  
The sun still moved.

\textbf{YOUNG ACOLYTE}  
It always moves—until it doesn’t.

\textbf{— WARRIORS’ QUARTERS}

Yacanex sits with Malinal and others from training. They clean weapons carefully, hands steadier than before.

\textbf{MALINAL}  
I heard three captives walked back under their own power.

\textbf{YACANEX}  
One corrected my words.

She looks at him.

\textbf{MALINAL}  
Corrected how?

\textbf{YACANEX}  
I said “I kill.”  
He told me what to say instead.

They sit with that.

\textbf{— SCRIBES’ ARCHIVE}

TLAMACAZQUI records numbers onto a fresh tablet. He pauses, erases a mark, rewrites it smaller.

\textbf{TLAMACAZQUI}  
(muttering)  
Unprecedented.

\textbf{— TLACAELEL’S HOME}

Tlacaelel enters quietly. Citlali waits, seated near a lamp.

\textbf{CITLALI}  
You’re late.

\textbf{TLACAELEL}  
The field held.

She studies his face.

\textbf{CITLALI}  
Is that good?

Tlacaelel considers.

\textbf{TLACAELEL}  
It means we didn’t collapse today.

Citlali pours him water.

\textbf{CITLALI}  
People are talking.

\textbf{TLACAELEL}  
They should.  
Silence would mean it failed.

He drinks, then sets the cup aside.

\textbf{TLACAELEL}  
There was one death.

Citlali nods, not surprised.

\textbf{CITLALI}  
There is always one.

\textbf{TLACAELEL}  
Yes.

He sits heavily.

\textbf{TLACAELEL}  
What troubles me is not the death.  
It’s that everyone noticed it.

Citlali sits beside him.

\textbf{CITLALI}  
You wanted them to count.

\textbf{TLACAELEL}  
I wanted them to understand.

Outside, distant drums begin—not victory, not mourning. Accounting.

\textbf{CAMERA:} Moves upward, out through the roof opening, over the city as lights emerge one by one.

The city remains.

\textbf{CUT TO:}

\section*{Scene X}

\textbf{INT. GREAT TEMPLE — INNER CHAMBER — NIGHT}

The chamber is enclosed and severe, lit by a ring of torches set low against the walls. Their flames throw long, wavering shadows that distort the painted glyphs of gods and stars.

\textbf{CAMERA:} Static at first, allowing the space to assert itself. Then a slow, deliberate push inward.

\textbf{TIZOC} stands with three \textbf{SENIOR PRIESTS}. Before them lies a bark-paper report sealed with the Emperor’s mark. It has already been opened, read, and reread.

\textbf{PRIEST ONE}  
Only one body.

\textbf{PRIEST TWO}  
Three captives returned alive.

\textbf{PRIEST THREE}  
The calendar did not object.

Tizoc says nothing. He traces a finger along a carved stone channel meant to carry blood during major rites.

\textbf{TIZOC}  
The danger is not that the sun stopped.  
The danger is that it did not notice.

The priests exchange uneasy glances.

\textbf{PRIEST ONE}  
The people speak carefully.  
They are not afraid.

\textbf{TIZOC}  
Fear is a language.  
If it is not spoken often enough,  
it is forgotten.

A pause.

\textbf{PRIEST TWO}  
What do we advise?

Tizoc turns slowly.

\textbf{TIZOC}  
We advise vigilance.  
We advise that restraint must still taste like blood.

He gestures toward the altar.

\textbf{TIZOC}  
Increase the ritual offerings between battles.  
Make absence visible.

The priests bow, understanding.

\textbf{CUT TO:}

\textbf{INT. MOCTEZUMA’S QUARTERS — SAME NIGHT}

Moctezuma sits alone now, listening to voices outside—warriors talking, not boasting, not lamenting.

\textbf{CAMERA:} Close on Moctezuma’s face as the sound filters in.

\textbf{WARRIOR (O.S.)}  
They stopped when the whistle sounded.

\textbf{ANOTHER WARRIOR (O.S.)}  
I didn’t think anyone would.

Moctezuma closes his eyes briefly.

\textbf{MOCTEZUMA}  
(quiet)  
Neither did I.

Cuauhtémoc enters, energized.

\textbf{CUAUHTÉMOC}  
They’re adapting.  
Fast.

Moctezuma looks at him sharply.

\textbf{MOCTEZUMA}  
Adaptation is not loyalty.

\textbf{CUAUHTÉMOC}  
No—but it is survival.

Moctezuma stands.

\textbf{MOCTEZUMA}  
That is what worries me.

\textbf{CUT TO:}

\textbf{INT. SCRIBES’ ARCHIVE — NIGHT}

TLAMACAZQUI and two younger scribes work by lamplight. Tablets and codices are stacked higher than usual.

\textbf{CAMERA:} Overhead, showing rows of recorded numbers extending beyond previous tallies.

\textbf{YOUNG SCRIBE}  
Where do we record captives who return?

TLAMACAZQUI hesitates.

\textbf{TLAMACAZQUI}  
We make a new column.

\textbf{YOUNG SCRIBE}  
And if it fills?

TLAMACAZQUI looks at the growing stacks.

\textbf{TLAMACAZQUI}  
Then the city will have to decide  
what kind of empire it is keeping track of.

He marks the tablet carefully.

\textbf{CUT TO:}

\textbf{EXT. CITY ROOFTOPS — NIGHT}

Xochitl runs across rooftops, carrying messages under moonlight. She pauses at a high point, looks out over the city.

Torches burn steadily. No fires. No alarms.

She fingers the death whistle at her throat but does not raise it.

\textbf{XOCHITL}  
(to herself)  
Not yet.

She runs on.

\textbf{CAMERA:} Pulls up above the city, the grid of canals and causeways visible like a living diagram—ordered, fragile, held together by agreement.

\textbf{CUT TO:}

\section*{Scene XI}

\textbf{INT. COUNCIL CHAMBER — DAY}

The chamber is fuller than before. Delegates from allied and rival city-states sit along the stone benches—distinct cloaks, different insignia, careful spacing. The atmosphere is no longer speculative. It is procedural.

\textbf{CAMERA:} Slow pan across faces unfamiliar to the audience, then settling on Tlacaelel at the center, now standing beside a low stone table covered in codices and marked tablets.

\textbf{TLACAELEL}  
The field at Xochitlan held.  
The field at Tepetzinco held.  
At each, captives exceeded fatalities by more than ten to one.

A murmur runs through the chamber—not awe, not fear. Calculation.

A \textbf{DELEGATE FROM CHALCO} leans forward.

\textbf{CHALCO DELEGATE}  
And the calendar?

\textbf{TLACAELEL}  
Uninterrupted.

\textbf{HUEXOTZINCO DELEGATE}  
You ask us to expose our warriors  
without the promise of conclusion.

Tlacaelel inclines his head.

\textbf{TLACAELEL}  
You expose them every season already.  
I am offering predictability.

Moctezuma sits among the captains, listening closely.

\textbf{MOCTEZUMA}  
Predictability favors those who can plan longest.

Tlacaelel meets his gaze.

\textbf{TLACAELEL}  
Which is why agreements must be mutual  
and violations immediate.

An elder from Texcoco speaks.

\textbf{TEXCOCAN ELDER}  
And who enforces these agreements?

Tlacaelel gestures outward, encompassing the room.

\textbf{TLACAELEL}  
Everyone who benefits from them.

A beat.

\textbf{CHALCO DELEGATE}  
That assumes restraint is valued more than advantage.

\textbf{TLACAELEL}  
It assumes advantage lasts longer  
when everyone expects to meet again.

Nezahualcoyotl watches the exchange, thoughtful.

\textbf{NEZAHUALCOYOTL}  
You are not abolishing war.  
You are extending memory.

The delegates fall silent, absorbing this.

\textbf{CUT TO:}

\textbf{EXT. MULTIPLE FIELDS — MONTAGE}

\textbf{CAMERA:} A controlled montage—brief, precise images.

— Warriors lining up on a new field, banners raised.  
— A death whistle sounding, echoed by another across the field.  
— Grapples, controlled strikes, captives walking upright.  
— Scribes recording names rather than body counts.  
— Messengers departing immediately after battles, spreading results.

\textbf{CUT TO:}

\textbf{INT. SCRIBES’ ARCHIVE — DAY}

The archive has expanded. New shelves line the walls.

TLAMACAZQUI supervises as younger scribes struggle to keep pace.

\textbf{YOUNG SCRIBE}  
We’re running out of symbols.

\textbf{TLAMACAZQUI}  
Then we invent more.

He gestures to a growing stack labeled with a new glyph: \textbf{RETURNED}.

\textbf{CUT TO:}

\textbf{INT. MOCTEZUMA’S QUARTERS — DAY}

Moctezuma reviews reports alone. The trophies on the walls feel heavier now, less triumphant.

Cuauhtémoc enters.

\textbf{CUAUHTÉMOC}  
The men speak of fields, not enemies.

Moctezuma looks up sharply.

\textbf{MOCTEZUMA}  
That language will cost us eventually.

\textbf{CUAUHTÉMOC}  
Or save us.

Moctezuma sets the report down.

\textbf{MOCTEZUMA}  
Empires survive by being feared.  
They endure by being needed.

Cuauhtémoc waits.

\textbf{MOCTEZUMA}  
Tlacaelel is trying to trade one for the other.

\textbf{CUT TO:}

\textbf{EXT. JUNGLE PATH — DAY}

Xochitl runs again, older now, faster still. Her route intersects others—an emerging network.

She passes another runner and exchanges a nod. No words needed.

\textbf{CAMERA:} Pulls back to reveal multiple paths crisscrossing the landscape, converging toward the city like veins.

\textbf{CUT TO:}

\textbf{INT. TLACAELEL’S HOME — NIGHT}

Tlacaelel and Citlali sit together in quiet. He studies a map marked with fields and dates.

\textbf{CITLALI}  
It’s working.

\textbf{TLACAELEL}  
Yes.

\textbf{CITLALI}  
You don’t sound relieved.

Tlacaelel traces a boundary line with his finger.

\textbf{TLACAELEL}  
Systems that work invite dependence.  
Dependence sharpens betrayal.

Citlali rests a hand on his arm.

\textbf{CITLALI}  
Then plan for that too.

Tlacaelel looks at the map, then at her.

\textbf{TLACAELEL}  
I am.

\textbf{CAMERA:} Holds on the map as new markings are added—order accumulating, pressure building.

\textbf{CUT TO:}

\section*{Scene XII}

\textbf{EXT. FLOWER WAR FIELD — LATE AFTERNOON}

The light is lower now, longer shadows stretching across the marked perimeter stones. The field is familiar—too familiar. Warriors on both sides settle into position with practiced ease.

\textbf{CAMERA:} Begins with a wide, orderly composition, then slowly tightens, introducing asymmetry—small deviations in stance, timing, breath.

Drums begin their count.

On the Mexica side, \textbf{TLATOANI} stands with a mixed unit of experienced fighters and newer recruits, including \textbf{MALINAL}. Across the field, Tlaxcalan ranks form—but one section is tighter, more aggressive.

\textbf{TLATOANI}  
(low)  
Something’s wrong.

\textbf{MALINAL}  
They’re leaning forward.

The death whistle sounds. Short. Sharp.

The lines advance.

Almost immediately, a Tlaxcalan warrior breaks cadence, surging ahead of his unit. His macuahuitl is unmodified—edges sharp along its full length.

\textbf{CAMERA:} Jumps closer. The rhythm fractures.

He strikes hard. Too hard.

A Mexica warrior goes down, blood pooling quickly.

A collective intake of breath ripples through both sides.

\textbf{TLATOANI}  
Hold!

Some obey. Others don’t hear—or don’t want to.

The field begins to tilt toward chaos.

\textbf{CAMERA:} Whip-pans between micro-decisions: a blade pulled back at the last instant; another that follows through.

Malinal engages the rogue Tlaxcalan, deflecting a lethal strike, then stepping inside his guard. She locks his arm, forces him down.

\textbf{MALINAL}  
Stop!

He struggles wildly.

\textbf{TLAXCALAN ROGUE}  
This is a lie!

Nearby, Tlatoani drags a wounded man clear of the fray.

\textbf{TLATOANI}  
Whistle! Now!

At the field’s edge, an official raises the death whistle but hesitates—eyes darting, unsure.

\textbf{CAMERA:} Close on the whistle, trembling in his grip.

Tlacaelel stands behind him, rigid.

\textbf{TLACAELEL}  
Sound it.

The whistle screams—long, sustained.

For a moment, no one stops.

Then—one by one—fighters disengage. Blades lower. Breathing replaces shouting.

The rogue Tlaxcalan is held fast, disarmed, furious.

Across the field, Tlaxcalan commanders shout orders, pulling their men back.

Silence falls heavily.

\textbf{CAMERA:} Wide shot. The field is intact, but stained. The boundary stones hold.

Tizoc steps forward, face unreadable.

\textbf{TIZOC}  
This is what hunger looks like.

Tlacaelel does not answer.

Tlaxcalan envoys cross the field cautiously.

\textbf{TLAXCALAN ENVOY}  
That warrior acted without sanction.

\textbf{TLACAELEL}  
Intent does not erase outcome.

He looks at the fallen Mexica warrior being carried away—still alive, barely.

\textbf{TLACAELEL}  
This is your warning.  
The next violation ends the agreement.

The envoy nods, grim.

\textbf{TLAXCALAN ENVOY}  
Understood.

The rogue warrior is taken away—not executed, not forgiven.

\textbf{CAMERA:} Lingers on Malinal, hands shaking now that the danger has passed.

Tlatoani meets her eyes.

\textbf{TLATOANI}  
You held.

She nods once, swallowing hard.

\textbf{CAMERA:} Pulls back slowly, the field returning to stillness, but the ease is gone.

\textbf{CUT TO:}

\section*{Scene XIII}

\textbf{INT. IMPERIAL COUNCIL ANNEX — NIGHT}

This chamber is smaller than the main council hall, tighter, more functional. The walls are bare stone. Lamps burn low. This is where decisions are refined after speeches end.

\textbf{CAMERA:} Close, compressed framing. Faces feel nearer. The room feels heavier.

Present are \textbf{TLACAELEL}, \textbf{ITZCOATL}, \textbf{TIZOC}, \textbf{MOCTEZUMA}, and a small circle of senior nobles and war captains. No attendants. No scribes.

Silence holds for several breaths.

\textbf{ITZCOATL}  
One violation.

\textbf{TLACAELEL}  
One too many.

\textbf{MOCTEZUMA}  
Or one too few.

Tlacaelel turns to him.

\textbf{TLACAELEL}  
You wanted proof the system would break.

\textbf{MOCTEZUMA}  
I wanted proof it could survive contact with instinct.

Tizoc steps forward.

\textbf{TIZOC}  
What you witnessed today was appetite remembering itself.

\textbf{TLACAELEL}  
And restraint reasserting itself.

\textbf{TIZOC}  
Barely.

Itzcoatl raises a hand.

\textbf{ITZCOATL}  
Enough.

He looks at Tlacaelel.

\textbf{ITZCOATL}  
Your design assumes mutual fear of consequence.  
Fear fades.

\textbf{TLACAELEL}  
Then consequence must be visible.

A noble frowns.

\textbf{NOBLE}  
You mean punishment?

\textbf{TLACAELEL}  
I mean accounting.

He steps toward the low table, places a tablet down.

\textbf{TLACAELEL}  
Every violation is recorded.  
Every city-state that breaks the field  
loses its next protection.

\textbf{MOCTEZUMA}  
You would invite retaliation.

\textbf{TLACAELEL}  
I would make it predictable.

Tizoc studies him closely.

\textbf{TIZOC}  
And sacrifice?

Tlacaelel meets his gaze.

\textbf{TLACAELEL}  
Violators provide it.

A stillness falls.

\textbf{ITZCOATL}  
You are binding blood to law.

\textbf{TLACAELEL}  
Blood already binds everything.  
I’m only deciding where it flows.

Moctezuma considers this, expression shifting—not agreement, not refusal.

\textbf{MOCTEZUMA}  
You’re turning mercy into leverage.

\textbf{TLACAELEL}  
Mercy without leverage is sentiment.

A captain speaks cautiously.

\textbf{CAPTAIN}  
And the warriors?

\textbf{TLACAELEL}  
They will know exactly  
what breaking restraint costs.

Tizoc exhales slowly.

\textbf{TIZOC}  
Then the gods will be fed  
by those who forget them.

Itzcoatl nods once.

\textbf{ITZCOATL}  
Proceed.

He looks at Tlacaelel.

\textbf{ITZCOATL}  
But understand this—  
you have made yourself necessary.

\textbf{TLACAELEL}  
I always was.

\textbf{CAMERA:} Holds on Moctezuma’s face—calculating, uneasy—as the others begin to disperse.

\textbf{CUT TO:}

\textbf{EXT. CITY WALLS — DAWN}

The city wakes. Runners depart in multiple directions. New banners are raised—subtle markings indicating fields now governed by stricter terms.

\textbf{CAMERA:} High aerial view of Tenochtitlan as a web of obligation tightens around it.

The city holds.

For now.

\textbf{CUT TO:}

\section*{Scene XIV}

\textbf{EXT. JUNGLE RIDGE — DAY}

The ridge overlooks a narrow valley where two city-state banners face one another across a newly designated field. The markers are fresh. The ground has not yet learned the weight of bodies.

\textbf{CAMERA:} Long lens from elevation, compressing distance so the opposing sides appear closer than they are—potential energy without release.

\textbf{XOCHITL}, older now, harder in the face, arrives at a run. She does not slow until she reaches a small cluster of officials from both sides. Her whistle hangs heavy at her chest.

A Tlaxcalan \textbf{FIELD CAPTAIN} studies her, wary.

\textbf{FIELD CAPTAIN}  
You are late.

\textbf{XOCHITL}  
No.  
You are early.

She hands him a sealed tablet. He breaks it, reads, stiffens.

\textbf{FIELD CAPTAIN}  
These terms are narrower.

\textbf{XOCHITL}  
They always are after blood spills  
where it shouldn’t.

A Mexica \textbf{FIELD SCRIBE} steps forward.

\textbf{FIELD SCRIBE}  
The violation has been recorded.  
This field is now bound by secondary sanction.

The Tlaxcalan captain looks between them.

\textbf{FIELD CAPTAIN}  
You are no longer inviting us to fight.

Xochitl meets his gaze evenly.

\textbf{XOCHITL}  
No.  
We are scheduling it.

The captain exhales sharply.

\textbf{FIELD CAPTAIN}  
And if we refuse?

Xochitl touches the whistle—not raising it, only acknowledging its presence.

\textbf{XOCHITL}  
Then the next field will not be marked.  
And what happens there  
will not be counted.

Silence stretches.

The captain folds the tablet carefully.

\textbf{FIELD CAPTAIN}  
Your architect is turning war into stone.

\textbf{XOCHITL}  
Stone lasts longer than anger.

\textbf{CUT TO:}

\textbf{EXT. SAME FIELD — LATER}

Warriors assemble. The atmosphere is different—less curiosity, more caution.

\textbf{CAMERA:} Close on faces. No excitement. Calculation.

The whistle sounds.

Combat begins—controlled, restrained—but now threaded with tension. No one forgets the sanction.

\textbf{CUT TO:}

\textbf{INT. SCRIBES’ ARCHIVE — DAY}

New tablets are added under a heading marked with a darker glyph: \textbf{VIOLATION HISTORY}.

TLAMACAZQUI watches younger scribes hesitate before making entries.

\textbf{YOUNG SCRIBE}  
Once it’s written here—

\textbf{TLAMACAZQUI}  
—it doesn’t fade.

\textbf{CUT TO:}

\textbf{INT. MOCTEZUMA’S QUARTERS — NIGHT}

Moctezuma listens as a captain reports.

\textbf{CAPTAIN}  
The runners carry authority now.  
Men stop when they see them coming.

Moctezuma nods slowly.

\textbf{MOCTEZUMA}  
Then the battlefield has moved.

\textbf{CUT TO:}

\textbf{INT. TLACAELEL’S HOME — NIGHT}

Tlacaelel studies a growing network of fields and routes scratched into a wooden board. Citlali watches from the doorway.

\textbf{CITLALI}  
You’ve made yourself unavoidable.

\textbf{TLACAELEL}  
That was always the risk.

\textbf{CITLALI}  
And when you’re gone?

Tlacaelel does not answer immediately.

\textbf{TLACAELEL}  
Then it will decide  
whether it was ever about me.

\textbf{CAMERA:} Holds on the board—lines, intersections, dates—an empire rendered as schedule and constraint.

\textbf{CUT TO:}

\section*{Scene XV}

\textbf{EXT. FLOWER WAR FIELD — MIDDAY}

The field is immaculate. Boundary stones are straightened. Banners hang at precise heights. Even the ground seems groomed, packed smooth by repetition rather than chaos.

\textbf{CAMERA:} A slow, almost reverent lateral movement across the field, emphasizing order—too much order.

Warriors assemble with practiced efficiency. Faces are calm. There is no anticipation, no fear. Only routine.

Tlacaelel stands with Tizoc and a small group of observers. Scribes wait with tablets already prepared.

\textbf{TIZOC}  
They know the steps now.

\textbf{TLACAELEL}  
Yes.

\textbf{TIZOC}  
The gods receive their due on schedule.  
The sun rises without argument.

Tlacaelel watches the warriors closely.

\textbf{TLACAELEL}  
Do you hear anything?

Tizoc listens. Drums beat steadily. Whistles wait.

\textbf{TIZOC}  
I hear compliance.

Tlacaelel nods, but uneasily.

The death whistle sounds. The battle begins.

\textbf{CAMERA:} Moves among the fighters. The choreography is precise, almost rehearsed. Captures occur quickly. Resistance is measured. Even struggle feels abbreviated.

Yacanex subdues an opponent with minimal effort. The enemy barely resists.

\textbf{YACANEX}  
(low, almost apologetic)  
Forgive me.

The opponent nods, resigned.

Nearby, Malinal executes a flawless capture. No blood is spilled.

The drums stop early.

The whistle sounds again.

It is over.

Scribes move in immediately, recording numbers before the dust settles.

\textbf{SCRIBE}  
Three captives.  
No fatalities.

A murmur of approval ripples through the officials.

Tlacaelel does not join it.

\textbf{TLACAELEL}  
Too clean.

Tizoc turns.

\textbf{TIZOC}  
You wanted fewer deaths.

\textbf{TLACAELEL}  
I wanted tension.  
This feels… anesthetized.

Tizoc considers.

\textbf{TIZOC}  
Ritual without danger becomes theater.

Tlacaelel watches as captives are escorted away with bureaucratic calm.

\textbf{TLACAELEL}  
And theater teaches the wrong lesson.

\textbf{CUT TO:}

\textbf{INT. COUNCIL ANNEX — LATER}

Delegates review reports. The mood is satisfied.

\textbf{DELEGATE}  
Stability has been achieved.

\textbf{ANOTHER}  
Predictability benefits everyone.

Tlacaelel stands apart, arms folded.

\textbf{TLACAELEL}  
Predictability also benefits those  
who want to break it cleanly.

A delegate waves this off.

\textbf{DELEGATE}  
You worry too much.

\textbf{TLACAELEL}  
That is my role.

\textbf{CUT TO:}

\textbf{EXT. JUNGLE ROUTE — SUNSET}

Xochitl runs past a group of young runners being trained. She slows, watching them practice.

\textbf{CAMERA:} Close on her face—pride mixed with concern.

\textbf{XOCHITL}  
(to a trainee)  
Do not memorize routes.  
Learn to notice silence.

The trainee looks confused.

\textbf{TRAINEE}  
Silence?

Xochitl looks toward the distant horizon.

\textbf{XOCHITL}  
Yes.  
That’s where danger hides.

\textbf{CUT TO:}

\textbf{INT. TLACAELEL’S HOME — NIGHT}

Tlacaelel sits alone with his codex. The pages are full now. No margins remain.

Citlali enters quietly.

\textbf{CITLALI}  
They say it’s perfect.

Tlacaelel closes the codex.

\textbf{TLACAELEL}  
Nothing that involves people ever is.

\textbf{CITLALI}  
What happens next?

Tlacaelel looks at the filled pages, then at her.

\textbf{TLACAELEL}  
Now we find out  
who was only waiting.

\textbf{CAMERA:} Holds on the closed codex—dense, complete, fragile.

\textbf{CUT TO:}

\section*{Scene XVI}

\textbf{EXT. COASTAL RIDGE — DAWN}

The horizon is pale, almost colorless. Mist hangs low over the water. Waves break with indifferent regularity.

\textbf{CAMERA:} Begins on the rhythmic motion of the sea, familiar and unthreatening, then tilts upward to reveal shapes that do not belong.

Three dark silhouettes sit offshore—angular, rigid, unmoving in a way that defies the water beneath them.

Sails.

They do not flutter like banners. They hold.

A group of \textbf{COAST WATCHERS} stands on the ridge, squinting, unsettled.

\textbf{WATCHER ONE}  
Are they islands?

\textbf{WATCHER TWO}  
Islands do not move.

The ships shift slightly, correcting against the current.

Silence deepens.

One watcher raises a runner’s horn, then hesitates.

\textbf{CUT TO:}

\textbf{INT. GREAT TEMPLE — ROOF — MORNING}

The city stretches below, orderly, awake. Tlacaelel stands with Tizoc and several senior priests. The air is thin here. The sky feels closer.

A \textbf{BREATHLESS RUNNER} arrives, bows deeply.

\textbf{RUNNER}  
From the coast.

He hands over a small carved tablet. Tizoc reads it first. His expression changes—not fear, not anger. Recognition.

\textbf{TIZOC}  
The calendar did not name this.

Tlacaelel takes the tablet, studies the glyphs.

\textbf{TLACAELEL}  
Because this does not arrive on cycles.

Below them, the city hums as usual.

\textbf{TIZOC}  
Do we sound the alarm?

Tizoc turns toward a ceremonial \textbf{DEATH WHISTLE}, larger than the runners’ versions, inlaid with jade and bone. It rests on a stone stand, unused for generations.

\textbf{TLACAELEL}  
Not an alarm.

\textbf{TIZOC}  
Then what?

Tlacaelel looks out toward the east, toward water he cannot see from here.

\textbf{TLACAELEL}  
A boundary marker.

Tizoc lifts the whistle, weighs it in his hands.

\textbf{TIZOC}  
This sound was meant to frighten men  
who understood fear.

He raises it to his lips.

\textbf{CAMERA:} Tight close-up. Breath drawn in.

The whistle screams—long, layered, inhuman. The sound rolls outward across the city, across the lake, toward the unseen coast.

People below stop. Heads turn upward. Conversation dies.

The sound fades.

A pause.

Tlacaelel listens—not to the city, but beyond it.

\textbf{TLACAELEL}  
They will not recognize the field.

\textbf{TIZOC}  
No.

\textbf{TLACAELEL}  
They will not agree to days,  
or limits,  
or return.

Tizoc lowers the whistle.

\textbf{TIZOC}  
Then your architecture has reached its edge.

Tlacaelel nods slowly.

\textbf{TLACAELEL}  
All architectures do.

\textbf{CUT TO:}

\textbf{EXT. JUNGLE PATH — DAY}

Xochitl runs hard. Harder than we have ever seen her. She passes relay stations without stopping, shouting messages as she goes.

\textbf{XOCHITL}  
Unknown force.  
No banners.  
No field.

Runners scatter in multiple directions.

\textbf{CUT TO:}

\textbf{INT. MOCTEZUMA’S QUARTERS — DAY}

Moctezuma receives the same report. He reads it once, then again.

\textbf{CUAUHTÉMOC}  
Enemies?

Moctezuma looks up.

\textbf{MOCTEZUMA}  
Something worse.

\textbf{CUAUHTÉMOC}  
Worse than war?

Moctezuma folds the tablet carefully.

\textbf{MOCTEZUMA}  
Worse than agreement.

\textbf{CUT TO:}

\textbf{EXT. FLOWER WAR FIELD — EMPTY — LATE AFTERNOON}

The field sits unused. Banners hang motionless. No warriors arrive.

\textbf{CAMERA:} Slow, lingering shot. Wind moves grass where bodies once stood.

The rules remain.

The participants do not.

\textbf{CUT TO:}

\section*{Scene XVII}

\textbf{INT. TLACAELEL’S HOME — NIGHT}

The house is quieter than it has been in years. No visitors. No attendants waiting for instruction. The lamps are turned low, their light contained rather than welcoming.

Rain taps softly against the courtyard stones.

\textbf{CAMERA:} Static, framed wide. Tlacaelel sits alone at first, cross-legged before the codex. He does not open it.

Citlali enters slowly. She carries no tray, no cup. She stops when she sees him still sitting where he was hours ago.

\textbf{CITLALI}  
You didn’t eat.

Tlacaelel does not answer immediately.

\textbf{TLACAELEL}  
I wasn’t hungry.

She sits opposite him, mirroring his posture. Between them, the closed codex feels heavier than when it was full.

\textbf{CITLALI}  
The city is quiet in a way I don’t recognize.

\textbf{TLACAELEL}  
Yes.

\textbf{CITLALI}  
Not fear.

\textbf{TLACAELEL}  
No.  
Suspension.

Rain grows slightly louder.

Citlali studies his face.

\textbf{CITLALI}  
You always said fear was manageable.

\textbf{TLACAELEL}  
Because fear looks at you.

\textbf{CITLALI}  
And this?

Tlacaelel exhales slowly.

\textbf{TLACAELEL}  
This doesn’t.

He finally opens the codex. The pages are dense with marks, glyphs, dates, tallies. No space remains unused.

\textbf{CAMERA:} Close on the pages as his fingers trace them without reading.

\textbf{TLACAELEL}  
Everything here assumes recognition.  
That everyone agrees the field exists.

Citlali places her hand on her abdomen, unconsciously.

\textbf{CITLALI}  
And if they don’t?

Tlacaelel closes the codex again.

\textbf{TLACAELEL}  
Then my work becomes a record,  
not a tool.

A long pause.

\textbf{CITLALI}  
You once told me war should return  
like breath.

\textbf{TLACAELEL}  
I assumed lungs.

She nods, accepting the metaphor without comment.

\textbf{CITLALI}  
Do you regret it?

Tlacaelel looks at her sharply.

\textbf{TLACAELEL}  
No.

He softens slightly.

\textbf{TLACAELEL}  
I regret believing success would protect it.

Citlali shifts closer.

\textbf{CITLALI}  
The child kicked today.

Tlacaelel freezes.

\textbf{TLACAELEL}  
Today?

\textbf{CITLALI}  
Yes.

She takes his hand, places it against her belly.

\textbf{CAMERA:} Tight on his hand as a faint movement presses back.

Tlacaelel closes his eyes.

\textbf{TLACAELEL}  
I built something meant to last  
longer than a life.

\textbf{CITLALI}  
And now?

\textbf{TLACAELEL}  
Now I want something that lasts  
longer than a collapse.

She leans her forehead against his.

\textbf{CITLALI}  
Then stop thinking like an architect.

\textbf{TLACAELEL}  
And think like what?

\textbf{CITLALI}  
Like a parent.

The rain slows.

\textbf{CUT TO:}

\textbf{EXT. CITY ROOFTOPS — SAME NIGHT}

Xochitl stands at the edge of a roof, looking east. The city spreads behind her—ordered, luminous, vulnerable.

She holds the death whistle in her hand, not raising it.

\textbf{XOCHITL}  
(quietly)  
You don’t hear us yet.

She closes her fingers around the whistle and turns back toward the city, already moving.

\textbf{CAMERA:} Pulls upward as she disappears into shadowed alleys, the city holding its breath beneath her.

\textbf{CUT TO:}

\section*{Scene XVIII}

\textbf{EXT. COASTAL BEACH — MORNING}

The shoreline is wide and pale, bordered by low dunes and scrub. Waves roll in with patient regularity. The air smells of salt and rot.

\textbf{CAMERA:} A distant, flattened perspective—heat shimmer, horizon pressed close—making depth hard to judge.

The ships are closer now. Not moving with tide or wind, but against them.

Wooden hulls grind into sand.

Figures disembark.

They are armored in unfamiliar ways: metal that catches light too cleanly, cloth stiffened with foreign seams. Animals follow—\textbf{HORSES}. Their size and motion unsettle the watchers immediately.

A group of \textbf{COASTAL SCOUTS} waits at a cautious distance. None advance. None flee.

One scout lifts a hand signal.

\textbf{CUT TO:}

\textbf{EXT. JUNGLE EDGE NEAR SHORE — CONTINUOUS}

Xochitl arrives with a relay runner. They stop short of the tree line.

\textbf{CAMERA:} Tight on Xochitl’s face as she takes in the scene—ships, animals, men who do not pause to assess the field.

\textbf{XOCHITL}  
No banners.

\textbf{RUNNER}  
No drums.

A Spanish \textbf{SOLDIER} laughs loudly, pointing at the scouts. He says something unintelligible.

Another soldier raises a metal tube—\textbf{A MUSKET}. He fires it into the air.

The sound is explosive, tearing through the ambient noise.

Several scouts flinch despite themselves.

\textbf{CAMERA:} Shakes slightly with the report. Birds scatter violently.

\textbf{XOCHITL}  
That was not a signal.

She lifts the death whistle instinctively, then stops.

\textbf{XOCHITL}  
They are not listening.

A mounted officer rides forward, horse stamping and snorting.

He gestures toward the scouts, smiling.

\textbf{SPANISH OFFICER}  
(in Spanish, untranslated)

The scouts do not understand the words, but the posture is unmistakable.

\textbf{CUT TO:}

\textbf{INT. IMPERIAL COUNCIL ANNEX — DAY}

The report lands hard.

Tlacaelel listens as Xochitl speaks—precise, unemotional, detailed.

\textbf{XOCHITL}  
They do not wait.  
They do not mirror.  
They do not ask.

\textbf{TLACAELEL}  
Do they retreat when faced?

\textbf{XOCHITL}  
No.

\textbf{MOCTEZUMA}  
Do they bleed?

Xochitl hesitates.

\textbf{XOCHITL}  
Yes.  
But they do not seem to notice.

A silence follows.

\textbf{TIZOC}  
Then they are not participants.

\textbf{TLACAELEL}  
No.

Tlacaelel stands, walks to the wall map. He studies the coastline.

\textbf{TLACAELEL}  
They move without regard for fields.  
Without season.  
Without reciprocity.

Moctezuma’s voice is low.

\textbf{MOCTEZUMA}  
Then they are invaders.

Tlacaelel does not contradict him.

\textbf{TLACAELEL}  
They are outside the system.

\textbf{TIZOC}  
Outside the gods’ patience as well.

Itzcoatl speaks for the first time.

\textbf{ITZCOATL}  
What does your architecture say  
about forces that do not agree?

Tlacaelel does not turn.

\textbf{TLACAELEL}  
It says containment fails.

He faces them now.

\textbf{TLACAELEL}  
We must assume annihilation  
as their default.

Moctezuma nods grimly.

\textbf{MOCTEZUMA}  
Then we answer in a language they understand.

\textbf{TLACAELEL}  
If we do, we abandon what we built.

\textbf{MOCTEZUMA}  
If we do not, we lose everything anyway.

A beat.

\textbf{TLACAELEL}  
Then the question is no longer  
how to preserve the system.

He gestures to the map, the city.

\textbf{TLACAELEL}  
But how to preserve the people.

\textbf{CUT TO:}

\textbf{EXT. FLOWER WAR FIELD — SUNSET}

The field lies empty again. Wind moves grass. One banner has fallen, its pole snapped.

\textbf{CAMERA:} Slow push forward across the boundary stones.

They no longer define anything.

\textbf{CUT TO BLACK.}

\section*{Scene XIX}

\textbf{INT. ARMORY HALL — NIGHT}

The armory is no longer quiet. Torches burn brighter than usual. Weapons are being removed from racks at a pace that feels urgent but not panicked.

\textbf{CAMERA:} Moves laterally past rows of macuahuitl—some modified, some old and fully edged—then pauses where the distinction breaks down. Warriors choose without ceremony.

Moctezuma stands at the center, issuing instructions to captains. His voice is controlled, clipped.

\textbf{MOCTEZUMA}  
No designated fields.  
No whistles.  
No counting.

A captain hesitates.

\textbf{CAPTAIN}  
And the captives?

Moctezuma meets his eyes.

\textbf{MOCTEZUMA}  
If they surrender, take them.  
If they don’t, end it.

The captain nods and moves on.

At the far edge of the hall, Tlacaelel watches. He does not interfere.

\textbf{CAMERA:} Holds on Tlacaelel’s face as warriors pass him—men and women he trained to stop, now preparing to finish.

Xochitl approaches him quietly.

\textbf{XOCHITL}  
They’re mobilizing the old routes.

\textbf{TLACAELEL}  
Yes.

\textbf{XOCHITL}  
Do you want me to carry anything different?

Tlacaelel thinks for a moment.

\textbf{TLACAELEL}  
Carry warnings.  
Not orders.

Xochitl nods once.

\textbf{XOCHITL}  
And if they don’t listen?

Tlacaelel’s answer is immediate.

\textbf{TLACAELEL}  
Then you run anyway.

She turns to go, then pauses.

\textbf{XOCHITL}  
You know this isn’t your failure.

Tlacaelel looks at her, almost surprised.

\textbf{TLACAELEL}  
It is still my responsibility.

She accepts that and leaves.

\textbf{CUT TO:}

\section*{Scene XX}

\textbf{EXT. JUNGLE APPROACH TO THE COAST — DAY}

Mexica forces move through dense growth. This is not ritual formation. Spacing is uneven. Scouts advance ahead, alert, tense.

\textbf{CAMERA:} Low and mobile, brushing leaves, obscured sightlines, sound-heavy—breath, insects, distant surf.

Tlatoani moves with Malinal and Yacanex. Their weapons are fully edged now.

\textbf{MALINAL}  
No markers.

\textbf{TLATOANI}  
No pauses either.

They hear a sharp metallic sound—armor shifting.

The first clash is sudden, chaotic.

Spanish soldiers fire muskets. Smoke erupts. Sound overwhelms space.

A Mexica warrior falls, stunned more than wounded. Another rushes forward and is cut down by steel.

\textbf{CAMERA:} Disoriented, partial frames. No choreography.

Malinal engages a mounted soldier. She tries to hook the leg as trained.

The horse tramples forward.

She is thrown aside hard.

Yacanex drags her clear.

\textbf{YACANEX}  
This isn’t a field.

\textbf{TLATOANI}  
No.

He charges anyway.

\textbf{CUT TO:}

\section*{Scene XXI}

\textbf{INT. GREAT TEMPLE — NIGHT}

The temple is crowded now. Priests chant louder, faster. Copal smoke thickens the air.

Tizoc stands before the altar, blood already staining the stone channel.

\textbf{CAMERA:} Static, frontal, ritual framing.

\textbf{TIZOC}  
The gods do not recognize restraint  
without threat.

A young priest hesitates beside him.

\textbf{YOUNG PRIEST}  
The offerings are increasing.

\textbf{TIZOC}  
Then so is the danger.

Tizoc pauses, then speaks more quietly.

\textbf{TIZOC}  
The Flower Wars taught us discipline.  
Now discipline must survive without them.

\textbf{CUT TO:}

\section*{Scene XXII}

\textbf{INT. TLACAELEL’S HOME — NIGHT}

Tlacaelel packs nothing. He sits with Citlali. The rain has returned.

\textbf{CITLALI}  
You could leave.

\textbf{TLACAELEL}  
And go where the field still exists?

She nods, knowing the answer already.

\textbf{CITLALI}  
What will you do?

\textbf{TLACAELEL}  
What I should have planned for earlier.

He places the codex into a cloth wrap—not to use, but to preserve.

\textbf{TLACAELEL}  
I will teach people how to remember  
when the structure fails.

Citlali rests her hand on his chest.

\textbf{CITLALI}  
Then teach our child first.

Tlacaelel holds her hand over the quiet movement beneath.

\textbf{CUT TO:}

\section*{Scene XXIII}

\textbf{EXT. CITY EDGE — DAWN}

Refugees move inward. Runners move outward. The order is gone, but motion remains.

\textbf{CAMERA:} High angle, tracking flows crossing unpredictably.

Xochitl runs past a fallen Flower War banner, picks it up, folds it carefully, and keeps running.

\textbf{XOCHITL}  
(to herself)  
Memory still moves.

She disappears into the city.

\textbf{CUT TO:}

\section*{Scene XXIV}

\textbf{EXT. OUTER DISTRICTS OF TENOCHTITLAN — DAY}

The city has changed shape. Barricades interrupt streets once open. Canals are crowded with canoes moving supplies rather than people. Smoke rises in places it never did before.

\textbf{CAMERA:} Slow tracking shot following Tlacaelel as he walks with effort now, older, leaning on a staff. He is unmarked, almost anonymous.

A small \textbf{GROUP OF YOUTHS} follows him—boys and girls, mixed ages. They are not warriors yet. They carry no weapons.

Tlacaelel stops beside a cleared patch of ground. Stones lie scattered.

\textbf{TLACAELEL}  
Sit.

They do.

\textbf{TLACAELEL}  
This was a field once.

A girl looks around, confused.

\textbf{GIRL}  
There are no markers.

Tlacaelel nods.

\textbf{TLACAELEL}  
Markers can be taken.  
Memory cannot.

He draws lines in the dirt with his staff—simple, deliberate.

\textbf{TLACAELEL}  
War teaches speed.  
Speed teaches mistakes.  
Restraint teaches recognition.

A boy raises a hand.

\textbf{BOY}  
Did it work?

Tlacaelel considers the question seriously.

\textbf{TLACAELEL}  
For a while.

The answer satisfies no one and everyone.

\textbf{TLACAELEL}  
Long enough to prove something important.

\textbf{GIRL}  
What?

Tlacaelel looks at the city behind them, then the open land beyond.

\textbf{TLACAELEL}  
That violence can remember itself.  
That humans can choose limits—  
even when the gods demand more.

The youths absorb this quietly.

\textbf{CUT TO:}

\section*{Scene XXV}

\textbf{EXT. JUNGLE RELAY STATION — DUSK}

The station is partially damaged but still standing. The glyphs are faded.

Xochitl arrives, older now, scarred, breath steady. She removes the death whistle from her neck and places it on the stone ledge.

\textbf{CAMERA:} Close on the whistle—worn smooth by decades of handling.

A young runner approaches.

\textbf{YOUNG RUNNER}  
Are we still running messages?

Xochitl nods.

\textbf{XOCHITL}  
Yes.

\textbf{YOUNG RUNNER}  
For whom?

Xochitl looks out into the darkening jungle.

\textbf{XOCHITL}  
For anyone who can still listen.

She lifts the whistle, hesitates, then blows—not the piercing scream of war, but a softer, controlled breath that produces a muted, mournful sound.

Not command.

Signal.

\textbf{CUT TO:}

\section*{Scene XXVI}

\textbf{INT. SCRIBES’ ARCHIVE — NIGHT}

The archive is smaller now. Many shelves are empty.

Tlacaelel sits with TLAMACAZQUI, both older. Between them lies the codex of the Flower Wars.

\textbf{CAMERA:} Static, respectful.

\textbf{TLAMACAZQUI}  
They will argue about this.

\textbf{TLACAELEL}  
Good.

\textbf{TLAMACAZQUI}  
They will say it failed.

\textbf{TLACAELEL}  
It ended.  
Those are not the same.

Tlacaelel closes the codex carefully.

\textbf{TLACAELEL}  
Write this clearly.

\textbf{TLAMACAZQUI}  
What?

\textbf{TLACAELEL}  
That for a moment,  
war was made to stop itself.

Tlamacazqui nods, begins writing.

\textbf{CUT TO:}

\section*{Scene XXVII}

\textbf{EXT. TOLTEC RUINS — SUNSET}

The same ruins as before. More eroded now. Eternal.

Tlacaelel stands alone. He touches the carved glyphs again—figures locked in controlled struggle.

\textbf{CAMERA:} Mirrors the opening scene’s composition.

\textbf{TLACAELEL}  
(softly)  
We remembered you.

Wind moves through broken stone.

A distant sound carries faintly—not a scream, not a drum.

A whistle.

\textbf{CUT TO:}

\section*{Scene XXVIII}

\textbf{EXT. FLOWER WAR FIELD — NIGHT}

The field is empty. Grass has grown tall. Boundary stones lie half-buried.

\textbf{CAMERA:} Very slow pullback under moonlight.

From far away, the sound of the death whistle rises—three short breaths, one long.

Not a call to fight.

A call to remember.

The sound fades.

The field remains.

\textbf{FADE TO BLACK.}

\bigskip

\begin{center}
\textit{End.}
\end{center}

\newpage
\section*{Historical Note}

This screenplay is inspired by historical figures, institutions, and practices of the Mexica (Aztec) civilization of the fifteenth century, but it is not a documentary reconstruction. It is a work of historical fiction that treats history as an intellectual and ethical landscape rather than a closed record of events.

Tlacaelel was a real historical figure: a high-ranking statesman, advisor, and reformer who played a decisive role in the political, religious, and military consolidation of the Mexica empire under rulers such as Itzcoatl and Moctezuma I. Many historians credit Tlacaelel with reshaping Mexica ideology, including the elevation of Huitzilopochtli, the rewriting of dynastic histories, and the institutionalization of practices that linked warfare, ritual, and cosmic maintenance.

The so-called “Flower Wars” (\textit{xōchiyāōyōtl}) are attested in multiple colonial-era sources. They are generally understood as ritualized conflicts—especially between the Mexica and their Tlaxcalan rivals—intended to capture prisoners for sacrifice, train warriors, and fulfill cosmological obligations. Scholarly debate continues regarding their precise origins, frequency, and political function. This screenplay adopts the interpretation that Flower Wars represented a constrained form of violence—war bounded by rules, calendars, and mutual recognition—rather than mere cruelty or spectacle.

Certain elements in the screenplay are speculative or dramatized. The degree to which Flower Wars were consciously designed as a systemic alternative to annihilatory warfare is not definitively known. The portrayal of Tlacaelel as an explicit architect of restraint reflects a thematic interpretation rather than a verified historical claim. Similarly, individual characters such as Citlali and Xochitl are fictional, created to give emotional, generational, and social texture to historical forces that are otherwise abstract in the sources.

Nahuatl terms, cosmological concepts, and ritual practices are used with care but are necessarily simplified for dramatic clarity. Human sacrifice is presented not as gratuitous violence but as it was understood within Mexica cosmology: a reciprocal obligation between humans and gods believed necessary to sustain the movement of the sun and the continuation of the world. This belief system is neither endorsed nor condemned by the film; it is treated as an internally coherent worldview whose logic shaped real historical decisions.

The arrival of the Spanish is depicted as a rupture rather than a conventional antagonist. Historically, Hernán Cortés and his allies exploited existing political tensions, disease, and asymmetries in military technology. This screenplay emphasizes a different asymmetry: the encounter between a system of ritualized, reciprocal warfare and an opponent who did not recognize the field itself. The collapse of the Flower Wars is therefore framed not as moral failure, but as a structural mismatch between incompatible conceptions of conflict.

Finally, this film resists the temptation to read history backward. The Mexica are not portrayed as inevitably doomed, nor are their institutions treated as primitive precursors to modern states. Instead, the Flower Wars are presented as a serious, if ultimately fragile, attempt to place limits on organized violence. The tragedy explored here is not that such limits failed forever, but that they existed at all—and were forgotten.


\end{document}
